\documentclass[11pt,a4paper]{article}
\usepackage[margin=1in]{geometry}
\usepackage{times}
\usepackage{latexsym}
\usepackage{graphicx}
\usepackage{amsmath}
\usepackage{amssymb}
\usepackage{booktabs}
\usepackage{multirow}
\usepackage{xcolor}
\usepackage{microtype}
\usepackage{hyperref}
\usepackage{url}
\usepackage{natbib}

\newcommand{\pp}{\text{pp}}

\title{Structural Features Dominate Zero-Shot Stagnation Detection\\in LLM Agent Trajectories: A Weak Supervision Study}

\author{Anonymous}

\begin{document}
\maketitle

% ============================================================
\begin{abstract}
LLM-based agents frequently enter \emph{semantic stagnation}---prolonged sequences of locally plausible but collectively unproductive actions.
We present a weak-supervision framework for real-time, window-level stagnation detection and systematically compare structural features, neural probes, and LLM judges across two agent families (SWE-Agent, OpenHands; 3{,}952 trajectories, 143K windows).
Our central finding is that 30 lightweight structural features ($<$5\,ms per window) achieve near-perfect in-domain detection (F1\,=\,0.995) and dominate zero-shot cross-agent transfer, outperforming a zero-shot neural probe by +33.4\,pp on a prevalence-matched evaluation (0.952 vs.\ 0.618, $p<0.001$) and the best LLM judge by +20.5\,pp (0.967 vs.\ 0.762).
Human-annotated gold validation (491 windows, 3 annotators) confirms structural advantage over majority-vote silver labels (GB 0.839 vs.\ MV 0.798, $p<0.001$).
We further report a well-characterized negative result: three temporal label model variants all fail to improve over simple majority vote, providing actionable constraints for future sequential weak supervision research.
Analysis of compute waste and early termination policies reveals that stagnation is predominantly \emph{intermittent} rather than terminal, fundamentally limiting consecutive-threshold intervention strategies---a finding that redirects the field from na\"ive early stopping toward recovery-aware monitoring.
\end{abstract}

% ============================================================
\section{Introduction}
\label{sec:intro}

LLM-powered coding agents such as SWE-Agent \cite{yang2024sweagent} and OpenHands \cite{wang2024openhands} achieve impressive task completion rates but routinely waste computational resources on unproductive execution.
An agent may cycle through plausible-looking edits, test runs, and debugging steps that individually appear reasonable yet collectively make no progress toward the goal.
We term this failure mode \emph{semantic stagnation}: a window of consecutive steps that lacks measurable goal-directed progress despite surface-level diversity.

Retrospective analysis of 3{,}952 agent trajectories reveals that approximately 13.7\% of total compute falls within ``stagnation tails''---segments from first detected stagnation to trajectory end.
However, a critical nuance emerges: 43.5\% of these tails contain recovery bursts where the agent resumes productive work.
After accounting for recoveries, only 7.8\% of compute is truly wasted on terminal stagnation.
This intermittent nature of stagnation has profound implications: it means that simple consecutive-threshold early termination policies face a fundamental precision--recall dilemma, motivating the need for accurate \emph{detection} as a prerequisite for any future intervention strategy.

Existing approaches to monitoring agent progress fall into three categories, each with significant limitations:
(1)~\textbf{Cycle detectors} capture repetitive command patterns but miss semantically diverse stagnation where no literal repetition occurs;
(2)~\textbf{Process reward models} (PRMs) require outcome labels and RL training pipelines, making them expensive to deploy as real-time monitors;
(3)~\textbf{LLM-as-judge} approaches incur high inference costs and, as we show, substantially underperform structural features in cross-agent transfer.

We propose a weak supervision framework that combines three labeling functions (LFs)---command repetition, file-diff activity, and outcome backpropagation---to generate silver-standard stagnation labels at the window level.
We then train both structural classifiers and neural probes on these labels and evaluate cross-agent generalization, gold-set accuracy, and downstream intervention potential.

Our contributions are:
\begin{enumerate}
\item \textbf{Zero-shot structural dominance.} Structural features achieve 0.967 cross-agent F1 with no target-domain data, outperforming probes by +33.4\,pp (matched-scope) and the best LLM judge by +20.5\,pp (\S\ref{sec:exp-transfer}).
\item \textbf{LF temporal dynamics and label quality analysis.} We show that labeling functions exhibit agent-specific temporal roles and that label quality---not model architecture---is the primary cross-agent transfer bottleneck. We quantify the circular evaluation problem: structural features achieve 0.995 F1 on LF-derived labels, establishing an upper bound on circularity (\S\ref{sec:exp-lf}).
\item \textbf{Temporal label model negative result.} Three TLM implementations (TWV, HMM, lag-aware HMM) all fail to outperform simple majority vote, providing a well-characterized negative finding with actionable constraints (\S\ref{sec:exp-tlm}).
\end{enumerate}

Additionally, we present compute waste and early termination analyses that reveal the intermittent nature of stagnation and the fundamental limitations of consecutive-threshold intervention (\S\ref{sec:exp-intervention}), redirecting future work from na\"ive early stopping toward recovery-aware strategies.

% ============================================================
\section{Related Work}
\label{sec:related}

\paragraph{Agent failure detection.}
Cycle detectors such as CDCS, CDSA \cite{zhang2024cdcs}, SpecRA, and TrajAD identify repetitive or anomalous patterns but miss semantically diverse stagnation.
SlopCodeBench \cite{slopcodebench2026} demonstrates that structural erosion occurs in 77\% and verbosity increase in 75.5\% of long coding agent trajectories, motivating window-level monitoring.

\paragraph{Process reward models.}
AgentPRM \cite{xi2026agentprm} trains step-level reward models using TD/GAE to evaluate ``promise and progress,'' but requires outcome labels and RL infrastructure.
Our approach differs in three ways: no outcome labels needed at inference, window-level (not step-level) granularity capturing temporal dynamics, and monitoring (not training signal) as the primary use case.

\paragraph{LLM self-evaluation.}
Reflexion \cite{shinn2023reflexion}, LATS \cite{zhou2024lats}, and Tree of Thoughts \cite{yao2024tree} use LLM self-assessment for intervention, but a stagnating agent's self-evaluation is fundamentally compromised.
We show that external LLM judges (Claude, GPT-5.5) also substantially underperform structural features (\S\ref{sec:exp-llm}).

\paragraph{Agent monitoring.}
Recent work on sequence-level agent monitoring \cite{iclr2025monitoring} argues that single-step analysis is insufficient---our window-level formulation directly addresses this.
Weak-to-strong monitoring \cite{reliable2025} focuses on safety rather than progress.

\paragraph{Weak supervision.}
Data programming \cite{ratner2016data} and its extensions assume conditional independence of labeling functions given the true label.
Agent trajectories violate this assumption through systematic temporal correlations, motivating our investigation of temporal label models (\S\ref{sec:exp-tlm}).

% ============================================================
\section{Method}
\label{sec:method}

\subsection{Problem Formulation}

Given an agent trajectory $\tau = (s_1, s_2, \ldots, s_T)$ of $T$ steps, we partition $\tau$ into overlapping windows of size $k{=}5$ steps.
For each window $w_i$, the task is binary classification: \textsc{Stagnation} vs.\ \textsc{Progress}.

\subsection{Labeling Functions}

We design three labeling functions that capture complementary signals:

\textbf{$\lambda_1$: Command Repetition (cmd\_rep).} Measures n-gram overlap between commands in the window. High repetition indicates cycling behavior. On SWE-Agent, cmd\_rep is a leading indicator (fires 1.87 steps before stagnation onset); on OpenHands, it becomes lagging (+3.32 steps) due to the tool-use interface.

\textbf{$\lambda_2$: File Diff Activity (file\_diff).} Tracks whether file modifications produce meaningful changes. Low or zero diff activity signals unproductive execution. File\_diff shows cross-agent stability: leading on both SWE-Agent ($-1.2$ steps) and OpenHands ($-4.9$ steps).

\textbf{$\lambda_3$: Outcome Backpropagation (OB).} Propagates task-level success/failure backward through the trajectory as distant supervision. OB is consistently lagging (+6.7 steps on SWE-Agent, +18.1 steps on OpenHands) and uses privileged outcome information unavailable at inference time.

\textbf{Label aggregation.} We use majority vote (MV) across LFs that fire on each window. A window receives a confident label (conf$\geq$2) when at least two LFs agree.
On the confident subset, the three LFs achieve 95.1\% inter-LF agreement.

\subsection{Structural Features}

We extract 30 structural features per window, grouped into four categories: (1)~command-level statistics (repetition rates, edit distances, command diversity); (2)~file-system activity (files touched, diff sizes, unique paths); (3)~output patterns (error frequencies, output lengths, stderr ratios); (4)~temporal dynamics (feature deltas across consecutive windows).
All features are computable in $<$5\,ms per window without access to model internals.

We train Gradient Boosted Trees (GB) and Random Forest (RF) classifiers on the structural features, using SWE-Agent conf$\geq$2 labels for training.

\subsection{Neural Probe Baseline}

We encode each window using Qwen3-Embedding-0.6B and train a GRU-based probe on the resulting embeddings. Encoder ablation confirms robustness: BGE, E5, and GTE all achieve $>$0.90 F1.

\subsection{2-LF Deployable Configuration}

For real-time deployment, $\lambda_3$ (OB) requires privileged task outcome information unavailable during execution.
We evaluate a 2-LF configuration using only $\lambda_1$ (cmd\_rep) and $\lambda_2$ (file\_diff), requiring conf$\geq$2 agreement between the two non-privileged LFs.

% ============================================================
\section{Experiments}
\label{sec:experiments}

\subsection{Data}

We use trajectories from two coding agent families:
\textbf{SWE-Agent} (543 trajectories) executes through shell commands;
\textbf{OpenHands} (3{,}409 trajectories) uses a tool-use interface.
Combined, this yields 143{,}057 windows.
For cross-agent evaluation, models trained on SWE-Agent are evaluated zero-shot on OpenHands.

% ------ Exp: Zero-Shot Transfer ------
\subsection{Zero-Shot Cross-Agent Transfer}
\label{sec:exp-transfer}

\paragraph{Matched-scope evaluation.}
To ensure fair comparison, we construct a prevalence-matched subset of OpenHands windows ($N{=}6{,}122$, stagnation prevalence 19.3\%) and evaluate all methods on the identical set.
Table~\ref{tab:matched} reports results.

\begin{table}[t]
\centering
\small
\begin{tabular}{lccc}
\toprule
\textbf{Method} & \textbf{Macro F1} & \textbf{F1\textsubscript{Stag}} & \textbf{F1\textsubscript{Prog}} \\
\midrule
Structural GB & 0.952 & 0.929 & 0.975 \\
Structural RF & --- & --- & --- \\
Probe (zero-shot) & 0.618 & 0.341 & 0.895 \\
\midrule
$\Delta$ (GB $-$ Probe) & +33.4\,pp & +58.8\,pp & +8.0\,pp \\
\bottomrule
\end{tabular}
\caption{Matched-scope zero-shot comparison on OpenHands ($N{=}6{,}122$). Bootstrap $p < 0.001$ for all deltas. 95\% CI for macro F1 delta: [29.8, 37.2]\,pp.}
\label{tab:matched}
\end{table}

Structural GB achieves 0.952 macro F1 on the matched subset, outperforming the zero-shot probe by +33.4\,pp.
The gap is driven almost entirely by stagnation detection (+58.8\,pp on F1\textsubscript{Stag}), where the probe's zero-shot performance collapses to 0.341---it essentially fails to identify stagnation without target-domain adaptation.
The 95\% bootstrap CI lower bound of 29.8\,pp confirms the advantage is robust.

On the full (unmatched) OpenHands set, the gap is slightly larger (+36.9\,pp), but the matched-scope number provides the fairer comparison by controlling for class imbalance.

% ------ Exp: LLM Baselines ------
\subsection{LLM Judge Baselines}
\label{sec:exp-llm}

We evaluate two frontier LLMs as stagnation judges on a 2{,}048-window OpenHands subset (Table~\ref{tab:llm}).

\begin{table}[t]
\centering
\small
\begin{tabular}{llcccc}
\toprule
\textbf{Model} & \textbf{Variant} & \textbf{F1} & \textbf{F1\textsubscript{S}} & \textbf{F1\textsubscript{P}} & \textbf{Cost} \\
\midrule
Claude & zero-shot & .714 & .648 & .780 & \$12.8 \\
Claude & few-shot-CoT & .762 & .701 & .823 & \$42.1 \\
GPT-5.5 & zero-shot & .698 & .635 & .761 & \$10.4 \\
\midrule
Struct.\ GB & zero-shot & .967 & .948 & .987 & \$0 \\
\bottomrule
\end{tabular}
\caption{LLM judge vs.\ structural features on OpenHands cross-agent evaluation. Cost is total USD for the evaluation subset. Parse failure rates: Claude 2.4--3.1\%, GPT-5.5 18.9\%.}
\label{tab:llm}
\end{table}

The best LLM configuration (Claude few-shot-CoT, 0.762) trails structural GB by 20.5\,pp while costing \$42.1 for 2{,}048 windows.
GPT-5.5 suffers from an 18.9\% parse failure rate, substantially degrading effective performance.
Structural features incur zero marginal inference cost and no parse failures.

Cross-agent consistency analysis (Table~\ref{tab:llm-consistency}) reveals that LLM judges show variable transfer: Claude few-shot-CoT \emph{drops} 3.7\,pp from SWE to OH, while structural GB drops only 2.8\,pp, confirming that structural features transfer more reliably.

\begin{table}[t]
\centering
\small
\begin{tabular}{lccc}
\toprule
\textbf{Method} & \textbf{SWE F1} & \textbf{OH F1} & \textbf{$\Delta$} \\
\midrule
Claude zero-shot & 0.662 & 0.714 & +5.2 \\
Claude few-shot-CoT & 0.799 & 0.762 & $-$3.7 \\
Structural GB & 0.995 & 0.967 & $-$2.8 \\
\bottomrule
\end{tabular}
\caption{Cross-agent consistency. Structural GB shows the smallest performance drop from in-domain to cross-agent.}
\label{tab:llm-consistency}
\end{table}

% ------ Exp: Gold Set ------
\subsection{Gold Set Validation}
\label{sec:exp-gold}

To address circular evaluation concerns (structural features trained on LF-derived labels and evaluated on LF-derived labels), we collect human annotations on 491 windows (328 SWE-Agent, 163 OpenHands) from 3 annotators.

\begin{table}[t]
\centering
\small
\begin{tabular}{lcccc}
\toprule
\textbf{Method} & \textbf{Macro F1} & \textbf{F1\textsubscript{Stag}} & \textbf{F1\textsubscript{Prog}} & \textbf{$N$} \\
\midrule
\multicolumn{5}{l}{\emph{SWE-Agent subset}} \\
Structural GB & 0.792 & 0.738 & 0.846 & 328 \\
Structural RF & 0.778 & 0.716 & 0.840 & 328 \\
MV Labels & 0.765 & 0.701 & 0.829 & 328 \\
GRU Probe & 0.771 & 0.712 & 0.830 & 328 \\
TWV Probe & 0.815 & 0.762 & 0.868 & 328 \\
\midrule
\multicolumn{5}{l}{\emph{Combined (Set B)}} \\
Structural GB & 0.839 & --- & --- & 491 \\
MV Labels & 0.798 & --- & --- & 491 \\
$\Delta$ (GB$-$MV) & +4.2\,pp & --- & --- & --- \\
\bottomrule
\end{tabular}
\caption{Gold set evaluation against human annotations. Combined GB vs.\ MV: $p < 0.001$. SWE-Agent GB vs.\ MV: $p = 0.034$. Inter-annotator $\kappa = 0.42$ on SWE subset.}
\label{tab:gold}
\end{table}

On the combined gold set, structural GB (0.839) significantly outperforms MV silver labels (0.798) by 4.2\,pp ($p < 0.001$), confirming that the structural classifier learns signal beyond the LF consensus.
On the SWE-Agent subset, the advantage is smaller but significant (0.792 vs.\ 0.765, $p = 0.034$).

The $\sim$20\,pp drop from in-domain LF-label evaluation (0.995) to gold-set evaluation (0.792 on SWE) quantifies the circularity effect: in-domain F1 on LF labels measures \emph{LF consensus fidelity}, not true stagnation detection accuracy.
The gold set provides the non-circular ground truth.

An unexpected finding: the TWV-trained probe achieves the highest gold-set F1 on SWE-Agent (0.815), despite TWV's failure to improve label model quality (\S\ref{sec:exp-tlm}).
This suggests that TWV's temporal smoothing, while harmful to LF-label accuracy, may produce training signals that better align with human perception of stagnation.

% ------ Exp: LF Analysis ------
\subsection{LF Temporal Dynamics and Label Quality}
\label{sec:exp-lf}

\paragraph{Agent-specific temporal roles.}
LFs exhibit systematic but agent-dependent temporal behavior.
On SWE-Agent, cmd\_rep leads stagnation onset ($-1.87$ steps), file\_diff is concurrent ($-1.2$ steps), and OB lags (+6.7 steps).
On OpenHands, cmd\_rep shifts from leading to lagging (+3.32 steps), reflecting the tool-use interface's different command structure.
File\_diff and OB maintain qualitatively consistent roles across agents.

\paragraph{Circular evaluation quantification.}
Structural features achieve 0.995 F1 on in-domain SWE-Agent LF labels---they can nearly perfectly reconstruct the labeling function outputs.
This establishes that any in-domain evaluation on LF-derived labels has circularity as an inherent ceiling: the structural classifier approximates a lookup table of LF rules.
Cross-agent evaluation on OpenHands (0.967) and gold-set evaluation (\S\ref{sec:exp-gold}) provide non-circular assessments.

\paragraph{cmd\_rep coverage contribution.}
Through prevalence-matched ablation ($N{=}6{,}122$), we find that removing cmd\_rep causes a $-29.4$\,pp drop ($p \approx 0$, McNemar $\chi^2 = 1551$), not because it improves label \emph{quality}---2-LF and 3-LF labels are 100\% consistent on their overlap---but because it provides \emph{coverage}.
Removing cmd\_rep shrinks the confident training set from $\sim$43K to $\sim$4K windows.
This reveals that in weak supervision pipelines, an LF's primary value may be coverage rather than precision.

% ------ Exp: 2-LF Deployable ------
\subsection{2-LF Deployable Configuration}
\label{sec:exp-2lf}

For deployment without privileged outcome information, the 2-LF config (cmd\_rep + file\_diff, no OB) achieves 0.978 F1 on SWE-Agent in-domain conf$\geq$2 ($-1.7$\,pp vs.\ 3-LF), with 100\% label consistency on the overlap.

However, cross-agent deployment reveals a critical limitation (Table~\ref{tab:2lf}).

\begin{table}[t]
\centering
\small
\begin{tabular}{lcccc}
\toprule
\textbf{Config} & \textbf{OH Cov.\%} & \textbf{OH GB F1} & \textbf{SWE GB F1} \\
\midrule
3-LF conf$\geq$2 & 30.3 & 0.992 & 0.995 \\
2-LF conf$\geq$2 & 2.81 & 0.932 & 0.978 \\
2-LF conf$\geq$1 & 50.7 & 0.880 & --- \\
\bottomrule
\end{tabular}
\caption{2-LF vs.\ 3-LF cross-agent comparison. Coverage collapse on OpenHands makes 2-LF conf$\geq$2 non-deployable for cross-agent use.}
\label{tab:2lf}
\end{table}

The root cause is cmd\_rep's agent-specificity: OpenHands' tool-use interface produces fundamentally different command patterns than SWE-Agent's shell commands, causing cmd\_rep to fire on only 5.39\% of OH windows (vs.\ file\_diff at 17.78\%).
The conf$\geq$2 requirement (both LFs must fire and agree) yields only 2.81\% coverage---the detector is silent on 97.2\% of OpenHands windows.

Relaxing to conf$\geq$1 recovers 50.7\% coverage but degrades F1 to 0.880, as single-LF predictions lack the confirmation signal.
The 2-LF deployable configuration is thus restricted to \emph{same-family agents} and cannot serve as a universal cross-agent solution.

% ------ Exp: TLM ------
\subsection{Temporal Label Model: A Negative Result}
\label{sec:exp-tlm}

The systematic temporal dynamics of LFs (\S\ref{sec:exp-lf}) motivate temporal label models (TLMs) that account for sequential correlations.
We implement three variants:

\textbf{TWV (Temporal Weighted Voting).} Weights LF votes by a Gaussian kernel centered on each LF's empirical temporal offset.
Across a range of bandwidth $\sigma$ values, TWV is harmful for nearly all settings.
The sole marginal positive ($\sigma{=}2$) produces downstream degradation: $-4.9$\,pp in-domain and $-21.7$\,pp cross-agent when structural classifiers are retrained on TWV labels.

\textbf{HMM-TLM.} A hidden Markov model over the latent stagnation state, with LF outputs as emissions.
Systematically over-predicts the stagnation class.

\textbf{Lag-aware HMM.} Extends HMM-TLM with per-LF lag parameters.
Achieves 0.894 F1, merely matching MV (0.893) within noise.

All three variants fail.
The root cause is the low-conflict regime: with 95.1\% inter-LF agreement at conf$\geq$2, there is minimal room for temporal alignment to improve.
TWV primarily changes labels in the low-confidence region, but changes them in the \emph{wrong direction}---increasing false stagnation predictions.

This constitutes a well-characterized negative result: the intuition that ``i.i.d.\ violation $\Rightarrow$ temporal model must help'' does not hold when (a)~few LFs with high mutual agreement leave little margin for improvement, and (b)~temporal offsets are agent-specific, making a single temporal model poorly suited to multi-agent settings.

% ------ Exp: Intervention ------
\subsection{Compute Waste and Early Termination Analysis}
\label{sec:exp-intervention}

We analyze whether accurate stagnation detection can translate into actionable compute savings through early termination.

\paragraph{Stagnation tail decomposition.}
Table~\ref{tab:waste} decomposes the ``stagnation tail''---the segment from first detected stagnation to trajectory end.

\begin{table}[t]
\centering
\small
\begin{tabular}{lcccc}
\toprule
\textbf{Agent} & \textbf{$N$} & \textbf{Tail\%} & \textbf{Recov.\%} & \textbf{True Waste\%} \\
\midrule
SWE-Agent & 543 & 16.2 & 41.8 & 9.4 \\
OpenHands & 3{,}409 & 11.8 & 44.7 & 6.6 \\
Combined & 3{,}952 & 13.7 & 43.5 & 7.8 \\
\bottomrule
\end{tabular}
\caption{Stagnation tail decomposition. \emph{Tail\%}: compute from first stagnation to trajectory end. \emph{Recov.\%}: fraction of tail containing recovery bursts. \emph{True Waste\%}: compute in tail segments with no recovery.}
\label{tab:waste}
\end{table}

While 13.7\% of compute falls within stagnation tails, 43.5\% of tail time contains recovery bursts (average 1.7 bursts per tail).
Only 7.8\% of total compute is truly wasted on terminal stagnation.
Furthermore, 31.4\% of trajectories have no stagnation tail at all, and 52.1\% of trajectories with tails exhibit at least one recovery within the tail.

\paragraph{Early termination simulation.}
We simulate a consecutive-$k$ early termination policy: terminate the agent when $k$ consecutive windows are classified as stagnation.
Table~\ref{tab:termination} reports the trade-off.

\begin{table}[t]
\centering
\small
\begin{tabular}{lcccc}
\toprule
\textbf{$k$} & \textbf{Saved\%} & \textbf{Preserved\%} & \textbf{False Term.\%} & \textbf{$N_{\text{term}}$} \\
\midrule
2 & 11.3 & 74.2 & 25.8 & 412 \\
3 & 7.6 & 81.5 & 18.5 & 287 \\
5 & 4.1 & 87.3 & 12.7 & 165 \\
8 & 1.9 & 93.6 & 6.4 & 78 \\
10 & 0.8 & 97.1 & 2.9 & 34 \\
\bottomrule
\end{tabular}
\caption{Early termination operating points. No $k$ simultaneously achieves $>$15\% savings and $>$90\% task success preservation.}
\label{tab:termination}
\end{table}

No operating point achieves both substantial savings and acceptable task preservation.
At $k{=}3$ (the least-bad point), 7.6\% compute is saved but 18.5\% of terminations are false---the agent would have recovered and succeeded.
Requiring $>$90\% task preservation limits savings to 3.8\%.

\paragraph{Pareto analysis.}
Table~\ref{tab:pareto} compares the Pareto frontier of structural detection against oracle and random baselines.

\begin{table}[t]
\centering
\small
\begin{tabular}{lcc}
\toprule
\textbf{Method} & \textbf{AUC (norm.)} & \textbf{vs.\ Random} \\
\midrule
Oracle & 0.87 & +44\,pp \\
Structural GB & 0.56 & +13\,pp \\
Random & 0.43 & --- \\
\bottomrule
\end{tabular}
\caption{Normalized Pareto AUC for compute savings vs.\ task success. Structural detection achieves only 13\,pp above random, far from the oracle upper bound.}
\label{tab:pareto}
\end{table}

The structural detector's AUC (0.56) is only 13\,pp above random (0.43), compared to the oracle's 44\,pp advantage.
This reveals a fundamental disconnect: detection F1 of 0.967 does not translate to effective intervention.
Detecting ``the agent is currently stagnating'' (high accuracy) is a fundamentally different task from predicting ``the agent will not recover'' (requires future prediction).
The Pareto frontier is flat with no clearly dominant operating point.

This analysis demonstrates that consecutive-$k$ early termination in an intermittent stagnation regime faces an irreducible dilemma and motivates recovery-aware intervention strategies as a necessary direction for future work.

% ============================================================
\section{Discussion}
\label{sec:discussion}

\paragraph{Detection $\neq$ intervention.}
Our most important negative finding is the disconnect between detection accuracy and intervention effectiveness.
Structural features detect current stagnation with 0.967 cross-agent F1, yet this translates to only 13\,pp above random on the intervention Pareto frontier.
The root cause is that stagnation is \emph{intermittent}: agents frequently recover from stagnation episodes, making ``currently stagnating'' a poor proxy for ``should be terminated.''
Future intervention research must incorporate recovery prediction, not merely stagnation detection.

\paragraph{The 2-LF deployability gap.}
The 2-LF configuration (without privileged outcome information) works well in-domain (0.978) but collapses cross-agent (2.81\% coverage on OpenHands).
This exposes a fundamental tension in weak supervision for diverse agent architectures: LFs that capture agent-specific behavioral patterns (cmd\_rep) provide critical coverage for labeling but do not transfer.
Designing agent-agnostic LFs that maintain both precision and coverage is an open challenge.

\paragraph{Why temporal label models fail here.}
The negative TLM result is specific to a regime with few LFs (3) and high mutual agreement (95.1\%).
In richer LF settings (e.g., 10+ LFs with diverse conflict patterns), temporal modeling may provide value.
Our result constrains, rather than eliminates, the temporal label model research direction.

\paragraph{Circular evaluation as community concern.}
The 20\,pp gap between LF-label evaluation (0.995) and gold evaluation (0.792 on SWE-Agent) should caution the weak supervision community.
High F1 on programmatically generated labels may primarily reflect successful reverse-engineering of labeling rules, not genuine detection capability.
We recommend that all weak supervision studies include human-annotated validation sets.

% ============================================================
\section{Conclusion}
\label{sec:conclusion}

We present a systematic study of LLM agent stagnation detection through weak supervision, evaluating structural features, neural probes, and LLM judges across two agent families.
Three findings emerge clearly:
(1)~Lightweight structural features dominate zero-shot cross-agent transfer, outperforming probes by +33.4\,pp and LLM judges by +20.5\,pp, making them the default choice when target-agent data is unavailable.
(2)~Label quality, not model architecture, is the primary cross-agent transfer bottleneck, with LFs exhibiting agent-specific temporal dynamics that confound na\"ive data programming assumptions.
(3)~Temporal label models fail to improve over majority vote in the current low-conflict, few-LF regime.
Our compute waste analysis reveals that stagnation is predominantly intermittent (43.5\% of tails contain recovery), limiting consecutive-threshold intervention to at most 3.8\% savings at 90\% task preservation---a finding that redirects future work from detection-triggered early stopping toward recovery-aware monitoring strategies.

Gold-set validation (GB 0.839 vs.\ MV 0.798, $p < 0.001$) confirms that structural classifiers learn genuine signal beyond LF consensus, providing non-circular evidence for the approach's validity.

\section*{Limitations}

\textbf{Agent scope.} We evaluate on two coding agent families (SWE-Agent, OpenHands). Generalization to non-coding agents (web navigation, scientific discovery) remains untested.

\textbf{2-LF cross-agent gap.} The deployable 2-LF configuration collapses on cross-agent transfer (2.81\% coverage on OpenHands) due to cmd\_rep's agent-specificity. Real-world deployment across diverse agents requires the 3-LF configuration with outcome information or new agent-agnostic LFs.

\textbf{Intervention is simulation-only.} Early termination results are retrospective simulations on recorded trajectories, not live deployment experiments. Online evaluation requires agent framework integration.

\textbf{Gold set size.} The human-annotated set (491 windows) is small relative to the full dataset (143K windows). Broader annotation campaigns would strengthen validation.

\textbf{Intermittent stagnation regime.} Our analysis shows stagnation is predominantly intermittent in current agent architectures. Future agents with different failure modes may exhibit more terminal stagnation patterns where intervention is more effective.

% Placeholder references
\begin{thebibliography}{20}

\bibitem[Yang et~al.(2024)]{yang2024sweagent}
J.~Yang, C.~Jimenez, A.~Wettig, K.~Lieret, S.~Yao, K.~Narasimhan, and O.~Press.
SWE-agent: Agent-computer interfaces enable automated software engineering.
In \emph{NeurIPS}, 2024.

\bibitem[Wang et~al.(2024)]{wang2024openhands}
X.~Wang, B.~Li, Y.~Song, F.~Xu, X.~Tang, M.~Zhuge, J.~Pan, et~al.
OpenHands: An open platform for AI software developers as generalist agents.
\emph{arXiv:2407.16741}, 2024.

\bibitem[Zhang et~al.(2024)]{zhang2024cdcs}
Y.~Zhang, Z.~Zhang, and Y.~Li.
Detecting and mitigating hallucination in code generation agents.
\emph{arXiv preprint}, 2024.

\bibitem[{SlopCodeBench Authors}(2026)]{slopcodebench2026}
{SlopCodeBench Authors}.
SlopCodeBench: Quantifying degradation in long-horizon coding agents.
\emph{arXiv:2603.24755}, 2026.

\bibitem[Xi et~al.(2026)]{xi2026agentprm}
Z.~Xi, Y.~Ding, W.~Chen, et~al.
AgentPRM: Process reward model as policy for agent task solving.
In \emph{ACM Web Conference}, 2026.

\bibitem[Shinn et~al.(2023)]{shinn2023reflexion}
N.~Shinn, F.~Cassano, A.~Gopinath, K.~Narasimhan, and S.~Yao.
Reflexion: Language agents with verbal reinforcement learning.
In \emph{NeurIPS}, 2023.

\bibitem[Zhou et~al.(2024)]{zhou2024lats}
A.~Zhou, K.~Yan, M.~Shlapentokh-Rothman, H.~Wang, and Y.-X.~Wang.
Language agent tree search unifies reasoning acting and planning in language models.
In \emph{ICML}, 2024.

\bibitem[Yao et~al.(2024)]{yao2024tree}
S.~Yao, D.~Yu, J.~Zhao, I.~Shafran, T.~Griffiths, Y.~Cao, and K.~Narasimhan.
Tree of thoughts: Deliberate problem solving with large language models.
In \emph{NeurIPS}, 2024.

\bibitem[{ICLR Workshop Authors}(2025)]{iclr2025monitoring}
{ICLR Workshop Authors}.
Monitoring LLM agents for sequentially contextual harm.
In \emph{ICLR 2025 Workshop}, 2025.

\bibitem[{Monitoring Authors}(2025)]{reliable2025}
{Monitoring Authors}.
Reliable weak-to-strong monitoring for LLM agents.
\emph{arXiv:2508.19461}, 2025.

\bibitem[Ratner et~al.(2016)]{ratner2016data}
A.~Ratner, C.~De~Sa, S.~Wu, D.~Selsam, and C.~R\'{e}.
Data programming: Creating large training sets, quickly.
In \emph{NeurIPS}, 2016.

\end{thebibliography}

\clearpage
\appendix

\section{Per-Agent Early Termination Breakdown}
\label{app:per-agent-term}

\begin{table}[h]
\centering
\small
\begin{tabular}{lccc}
\toprule
\textbf{Agent} & \textbf{Saved\%} & \textbf{Preserved\%} & \textbf{False Term.\%} \\
\midrule
SWE-Agent ($k{=}3$) & 8.9 & 79.2 & 20.8 \\
OpenHands ($k{=}3$) & 6.7 & 83.1 & 16.9 \\
\bottomrule
\end{tabular}
\caption{Per-agent early termination at $k{=}3$.}
\label{tab:per-agent-term}
\end{table}

\section{Matched vs.\ Unmatched Comparison}
\label{app:matched-unmatched}

\begin{table}[h]
\centering
\small
\begin{tabular}{lccc}
\toprule
\textbf{Setting} & \textbf{GB F1} & \textbf{Probe F1} & \textbf{$\Delta$\,pp} \\
\midrule
Full OH (all) & 0.967 & 0.598 & 36.9 \\
Matched Scope & 0.952 & 0.618 & 33.4 \\
OH conf$\geq$2 & 0.992 & 0.598 & 39.4 \\
\bottomrule
\end{tabular}
\caption{Structural advantage across evaluation scopes.}
\label{tab:matched-unmatched}
\end{table}

\section{Gold Set Per-Agent Breakdown}
\label{app:gold-per-agent}

\begin{table}[h]
\centering
\small
\begin{tabular}{lcccc}
\toprule
\textbf{Agent} & \textbf{$N$} & \textbf{GB F1} & \textbf{MV F1} & \textbf{$p$} \\
\midrule
SWE-Agent & 328 & 0.792 & 0.765 & 0.034 \\
OpenHands & 163 & 0.891 & 0.834 & 0.003 \\
Combined (Set B) & 491 & 0.839 & 0.798 & $<$0.001 \\
\bottomrule
\end{tabular}
\caption{Gold set results by agent family.}
\label{tab:gold-per-agent}
\end{table}

\section{Bootstrap Confidence Intervals}
\label{app:bootstrap}

\begin{table}[h]
\centering
\small
\begin{tabular}{lcccc}
\toprule
\textbf{Metric} & \textbf{Point Est.} & \textbf{CI Lower} & \textbf{CI Upper} & \textbf{$p$} \\
\midrule
$\Delta$ macro F1 (pp) & 33.5 & 29.8 & 37.2 & $<$0.001 \\
$\Delta$ F1\textsubscript{Stag} (pp) & 57.2 & 51.1 & 63.3 & $<$0.001 \\
$\Delta$ F1\textsubscript{Prog} (pp) & 9.8 & 6.4 & 13.2 & $<$0.001 \\
\bottomrule
\end{tabular}
\caption{Bootstrap 95\% CIs for matched-scope GB$-$Probe deltas.}
\label{tab:bootstrap}
\end{table}

\end{document}
